% =========================================================
% فهرست مطالب، تصاویر و جداوّل
% =========================================================
\tableofcontents
\newpage
\listoffigures
\newpage
% =========================================================
% ۱. توضیح کلی شیوه‌نامه\LTRfootnote{protocol} UART و قالب قاب مورد استفاده
% =========================================================
\section{پروتکل \texorpdfstring{\lr{UART}}{UART} و قالب قاب مورد استفاده}
شیوه‌نامه ارتباط سریال ناهمگام\LTRfootnote{Universal Asynchronous Receiver-Transmitter (UART)} یک استاندارد سخت‌افزاری است که در آن فرستنده و گیرنده نبض\LTRfootnote{pulse} ساعت\LTRfootnote{Clock} مشترکی ندارند. به منظور برقراری ارتباط صحیح، هر دو سمت باید بر روی سرعت انتقال و ساختار داده توافق داشته باشند.

در این آزمایش، قالب قاب\LTRfootnote{Frame Format} ارتباطی به طول ۱۰ بیت طراحی شده است که مشخصات آن به شرح زیر می‌باشد:
\begin{itemize}
    \item \textbf{بیت شروع\LTRfootnote{Start Bit}:} یک بیت با مقدار منطقی صفر (\lr{Low}) که خط را از حالت بی‌کار (\lr{IDLE}) خارج کرده و گیرنده را برای دریافت داده بیدار می‌کند.
    \item \textbf{بیت‌های داده\LTRfootnote{Data Bits}:} محتوای اصلی به طول ۷ بیت که ارسال آن‌ها از کم‌ارزش‌ترین بیت\LTRfootnote{Least Significant Bit (LSB)} یعنی \lr{D0} آغاز می‌گردد.
    \item \textbf{بیت توازن\LTRfootnote{Parity Bit}:} یک بیت از نوع توازن زوج\LTRfootnote{Even Parity} که از طریق عملگر \lr{XOR} بر روی ۷ بیت داده محاسبه می‌شود تا خطایابی در مقصد امکان‌پذیر گردد.
    \item \textbf{بیت توقف\LTRfootnote{Stop Bit}:} یک بیت با مقدار منطقی یک (\lr{High}) که پایان قاب را اعلام کرده و خط را مجدداً به حالت بی‌کار بازمی‌گرداند.
\end{itemize}

% =========================================================
% ۲. توضیح ماشین حالت فرستنده و نقش سیگنال‌ها
% =========================================================
\section{توضیح ماشین حالت فرستنده و نقش سیگنال‌های \lr{tx} و \lr{busy}}
پیمانه\LTRfootnote{Module} فرستنده بر پایه یک ماشین حالت متناهی\LTRfootnote{Finite State Machine (FSM)} با ۵ وضعیت اصلی طراحی شده است. به منظور تطبیق سرعت فرستنده و گیرنده، یک شمارنده داخلی (\lr{tick\_cnt}) در تمام وضعیت‌ها وظیفه دارد تا ماندگاری در هر بیت را به اندازه پارامتر \lr{BIT\_TICKS} تضمین نماید.

\subsection{تحلیل وضعیت‌های ماشین حالت}
گذر حالت‌ها در این پیمانه به شرح زیر است:

